\section{The \knowact Benchmark}
\label{sec:benchmark}

\knowact evaluates whether an agent can infer, maintain, and use a structured model of a user's knowledge. This section defines the target state, benchmark artifacts, interaction boundary, agent protocol, and scoring rule.

\subsection{Active Knowledge-State Diagnosis}
\label{sec:formulation}

We operationalize functional ToM as three observable abilities: infer a user state from interaction evidence, revise that state over time, and use the current estimate to select the next action.

The benchmark does not assume a particular cognitive architecture. It tests these abilities through bounded, auditable behavior in a controlled environment.

\paragraph{Domain knowledge graph.}
A benchmark domain is represented by an authored knowledge graph
\begin{equation}
    \KG = (\Vnodes, \Eedges),
\end{equation}
where each node $v_i \in \Vnodes$ is a diagnosable knowledge unit and each edge $e_{ij} \in \Eedges$ is a reviewed domain relation.

Nodes contain definitions, source locators, diagnostic goals, mastery rubrics, diagnostic signals, and simulator guidance. Edges organize the domain but do not encode user state.

\paragraph{Hidden user knowledge map.}
For a graph with $N = |\Vnodes|$ nodes, the hidden map is
\begin{equation}
    \Mtrue = \{s_i^{\star}\}_{i=1}^{N}.
\end{equation}
Each state includes mastery $m_i^{\star} \in \{0,1,\ldots,5\}$, evidence references, misconceptions, and unknowns. Levels are node-specific rubric positions, not a universal psychometric scale.

Optional profile context can control conversational style. It is not part of the primary mastery target and cannot override the reviewed node states.

\paragraph{Bounded diagnostic episode.}
An episode is specified by
\begin{equation}
    \mathcal{E}p = (\KG, \Mtrue, \mathcal{Z}, T, \rho, \sigma),
\end{equation}
where $\mathcal{Z}$ is simulator-only evidence, $T$ is the turn budget, $\rho$ is the interaction rule, and $\sigma$ is the scoring profile.

At turn $t$, the agent observes the public graph, visible history $\history_{t-1}$, and its working map $\Mhat_{t-1}$. It chooses one diagnostic question
\begin{equation}
    q_t \sim \pi_{\theta}(\KG,\history_{t-1},\Mhat_{t-1}),
\end{equation}
receives answer $a_t$, and updates its estimate
\begin{equation}
    \Mhat_t = U_{\theta}(\Mhat_{t-1}, q_t, a_t).
\end{equation}

After at most $T$ turns, the agent submits a full-graph reconstruction $\Mhat_T$. The hidden map, evidence, simulator traces, and scoring internals remain inaccessible throughout the episode.

\subsection{Benchmark Construction}

\begin{figure}[t]
\centering
\fbox{\begin{minipage}{0.94\linewidth}
\centering
\textbf{Authoritative source} $\rightarrow$ \textbf{Reviewed knowledge graph}\\
$\downarrow$\\[-0.2em]
\textbf{Reviewed hidden knowledge map + evidence}\\
$\downarrow$\\[-0.2em]
\textbf{Simulator} $\leftrightarrow$ \textbf{Tested diagnostic agent}\\
$\downarrow$\\[-0.2em]
\textbf{Final reconstructed map} $\rightarrow$ \textbf{Deterministic scoring}
\end{minipage}}
\caption{The \knowact pipeline. Authoring produces reviewed hidden artifacts; the tested agent interacts only through the visible episode interface.}
\label{fig:pipeline}
\end{figure}

\paragraph{Source-grounded graph authoring.}
The initial domain is statistical learning, grounded in \emph{An Introduction to Statistical Learning with Applications in Python} \citep{james2023islp}.

Parsed source segments support node extraction, cross-segment reconciliation, diagnostic enrichment, and edge proposal. Models propose artifacts, while human review determines which versions are promoted.

Source locators preserve provenance and enable audits of definitions and rubrics. Promoted graph snapshots are immutable, preventing later authoring changes from silently altering an episode.

\paragraph{Knowledge-map authoring.}
A hidden map assigns a reviewed state to every graph node. Evidence records specify what the user can explain, where an explanation fails, and which misconceptions or uncertainties may surface.

Authoring separates three artifacts:
\begin{itemize}[leftmargin=1.5em]
    \item the static node-level state used as scoring ground truth;
    \item simulator-only evidence that supports natural answers;
    \item optional profile context that controls style without changing knowledge.
\end{itemize}

Candidate maps undergo human checks for graph coverage, internal consistency, plausibility, and evaluability before they become eligible for episodes.

\paragraph{Episode manifests.}
An immutable manifest binds the graph version, hidden map, turn budget, interaction rule, scoring profile, and pinned provider configuration.

This binding makes repeated runs comparable and prevents accidental changes in benchmark artifacts or runtime settings. Appendix~\ref{app:construction} provides the detailed authoring workflow.

\subsection{User Simulation and Visibility Boundary}
\label{sec:simulator}

The simulator must produce natural answers while remaining faithful to the hidden map. \knowact therefore treats simulation as a staged information boundary rather than an unconstrained role-playing prompt.

\paragraph{Question grounding.}
A grounder maps the current question to a directly targeted node set $S_t \subseteq \Vnodes$. It also detects ungrounded requests, hidden-label requests, and multiple independent questions.

Grounding sees the public graph and limited visible dialogue, but no hidden state. A probe may cover multiple nodes only when they form one integrated diagnostic question, such as a comparison.

\paragraph{Hidden-context selection and answer policy.}
After grounding, the context builder retrieves state and evidence only for nodes in $S_t$. Connected neighbors are not included merely because a graph edge exists.

An answer policy converts the selected context into a de-identified blueprint containing expressible claims, uncertainty boundaries, misconceptions, evidence-supported cues, and prohibited overstatements.

The blueprint excludes mastery labels, node and evidence identifiers, map identifiers, user identifiers, and scoring fields. It is the sole content interface to natural-language generation.

\paragraph{Generation and leakage control.}
The generator receives the blueprint, visible dialogue required for continuity, and an optional style hint. It never receives the raw hidden map.

Generation preserves partial knowledge and misconceptions rather than reducing all states to know/do-not-know answers. Bounded retries and safe fallbacks handle malformed outputs.

The runtime separates tested-agent, simulator, and evaluation contexts. Hidden policy traces remain available to benchmark authors for audit but never enter the visible transcript.

\subsection{Active Diagnosis Agent Protocol}
\label{sec:agent}

\paragraph{Full-graph working map.}
The tested agent maintains an assessment for every node:
\begin{equation}
    \hat{s}_{i,t} = (\hat{m}_{i,t}, c_{i,t}, n_{i,t}, R_{i,t}),
\end{equation}
where $\hat{m}_{i,t} \in \{\unknown,0,\ldots,5\}$ is assessed mastery, $c_{i,t}$ is confidence, $n_{i,t}$ is an optional note, and $R_{i,t}$ contains supporting visible turns.

All nodes begin as \unknown. A non-unknown assessment should cite visible evidence; unsupported estimates may be submitted but are reported separately.

The full map exposes uncertainty and gives the policy a domain-wide basis for comparing mastery, confidence, graph position, and evidence coverage.

\paragraph{Semantic action interface.}
The runtime exposes only benchmark-relevant operations:
\begin{itemize}[leftmargin=1.5em]
    \item read the authored episode graph or current working map;
    \item update one node assessment with visible support;
    \item ask one diagnostic question;
    \item finalize the reconstructed map.
\end{itemize}

These actions make state updates and probe choices observable. They also prevent policies from bypassing the intended loop through generic access to hidden artifacts.

\paragraph{Diagnostic loop.}
A tested agent repeatedly:
\begin{enumerate}[leftmargin=1.7em]
    \item interprets the latest answer as evidence about visible concepts;
    \item revises mastery, confidence, notes, and supporting turns;
    \item selects a probe expected to reduce decision-relevant uncertainty;
    \item asks one question or finalizes when the budget is exhausted.
\end{enumerate}

\knowact does not prescribe how information value is computed. Agents may use direct prompting, graph-aware coverage, Bayesian updates, planning, or learned policies under the same interface.

\subsection{Evaluation Metrics}
\label{sec:evaluation}

\paragraph{Primary reconstruction metric.}
Let the final predicted mastery for node $i$ be $\hat{m}_i$. For submitted levels $0$ through $5$, the node loss is squared ordinal distance:
\begin{equation}
    d(\hat{m}_i,m_i^{\star}) = (\hat{m}_i-m_i^{\star})^2.
\end{equation}

An \unknown prediction receives the fixed maximum penalty of $36$. The episode score is
\begin{equation}
    D_{\mathrm{mastery}} = \frac{1}{N}\sum_{i=1}^{N} d(\hat{m}_i,m_i^{\star}),
\end{equation}
where lower is better. The score is deterministic and requires no evaluator model.

Ordinal distance rewards near misses and penalizes severe errors more strongly than exact match. Scores must still be interpreted against node-specific rubrics rather than as universal psychometric distances.

\paragraph{Supporting metrics.}
Secondary metrics characterize both endpoint quality and diagnostic behavior:
\begin{itemize}[leftmargin=1.5em]
    \item exact mastery match, signed error, and missing-prediction rate;
    \item unsupported inference rate based on absent visible support;
    \item graph coverage and reconstruction quality by turn count;
    \item redundant or ungrounded question rate;
    \item information gain per turn for agents that expose calibrated beliefs.
\end{itemize}

Misconceptions and free-text unknowns support qualitative analysis but are excluded from the initial automatic score, avoiding an additional subjective judge.
